\documentclass[12pt]{article}
\usepackage[utf8]{inputenc}
\usepackage[T1]{fontenc}
\usepackage[spanish]{babel}
\usepackage{amsmath,amssymb,amsthm}
\usepackage{mathtools}
\usepackage{enumitem}
\usepackage{geometry}
\usepackage{siunitx}
\usepackage{tikz}
\usepackage{hyperref}
\geometry{margin=1in}

\newtheorem{exercise}{Ejercicio}
\newtheorem{solution}{Solución}
\setlist[exercise]{leftmargin=0pt,labelsep=1em}
\setlist[solution]{leftmargin=1.5em}

\begin{document}

\begin{center}
  {\LARGE \textbf{Ejercicios Avanzados de Ecuaciones y Proporciones}}\

\[0.5em]
  {\large Con soluciones detalladas}
\end{center}

\bigskip

\section*{Instrucciones}
\begin{itemize}
  \item Resuelve cada ejercicio mostrando todos los pasos.
  \item Para las proporciones, indica si es directa, inversa o compuesta y modela la relación algebraicamente.
  \item Usa simplificación algebraica y verifica soluciones sustituyendo en la ecuación original.
\end{itemize}

\section*{Ejercicios}

\begin{exercise}[Ecuaciones lineales y fraccionarias]
Resuelve las siguientes ecuaciones:
\begin{enumerate}[label=(\alph*)]
  \item \(3x - 7 = 2(x+5)\).
  \item \(\dfrac{2x+1}{x-1} = 3\).
  \item \(\dfrac{4}{x} + \dfrac{3}{x+2} = 1\).
\end{enumerate}
\end{exercise}

\begin{exercise}[Ecuaciones cuadráticas y con parámetros]
\begin{enumerate}[label=(\alph*)]
  \item Resuelve \(x^2 - 5x + 6 = 0\).
  \item Para qué valores de \(k\) la ecuación \(x^2 + kx + 1 = 0\) tiene raíces reales y distintas.
  \item Resuelve \(x^2 - (m+1)x + m = 0\) en función de \(m\).
\end{enumerate}
\end{exercise}

\begin{exercise}[Sistemas no lineales]
Resuelve el sistema


\[
\begin{cases}
x^2 + y^2 = 25,\

\[4pt]
x - y = 1.
\end{cases}
\]


\end{exercise}

\begin{exercise}[Ecuaciones exponenciales y logarítmicas]
\begin{enumerate}[label=(\alph*)]
  \item Resuelve \(2^{x+1} = 8\).
  \item Resuelve \(3^{2x-1} = 27\).
  \item Resuelve \(\log_2(x+3) - \log_2(x-1) = 2\).
\end{enumerate}
\end{exercise}

\begin{exercise}[Ecuaciones racionales complejas]
Resuelve


\[
\frac{x^2-4}{x^2-1} = \frac{3x-2}{x+1}.
\]


Indica las restricciones del dominio.
\end{exercise}

\begin{exercise}[Proporciones directas e inversas]
\begin{enumerate}[label=(\alph*)]
  \item Si \(y\) es directamente proporcional a \(x^2\) y \(y(2)=8\), encuentra \(y(5)\).
  \item Si \(t\) es inversamente proporcional a \(v\) y \(t(4)=3\), encuentra \(t(12)\).
  \item Problema compuesto: la intensidad \(I\) de una luz varía directamente con la potencia \(P\) e inversamente con el cuadrado de la distancia \(d\). Si \(I=50\) cuando \(P=100\) y \(d=2\), calcula \(I\) cuando \(P=150\) y \(d=3\).
\end{enumerate}
\end{exercise}

\begin{exercise}[Problemas aplicados y modelado]
\begin{enumerate}[label=(\alph*)]
  \item Un tanque se llena con dos mangueras: la primera llena en 3 horas y la segunda en 5 horas. ¿Cuánto tardan juntas? Modela y resuelve.
  \item En una mezcla, la relación masa de A a B debe ser \(3:5\). Si se necesitan 40 kg de mezcla, ¿cuántos kg de A y B?
  \item Un objeto cae con posición \(s(t)=100-4.9t^2\). Encuentra el tiempo en que la velocidad es \(-19.6\ \mathrm{m/s}\) y verifica la posición.
\end{enumerate}
\end{exercise}

\bigskip
\section*{Soluciones}

\begin{solution}[Ejercicio 1]
\textbf{(a)} \(3x-7=2(x+5)\Rightarrow 3x-7=2x+10\Rightarrow x=17.\)

\textbf{(b)} \(\dfrac{2x+1}{x-1}=3 \Rightarrow 2x+1=3(x-1)=3x-3 \Rightarrow -x=-4 \Rightarrow x=4.\) Verificar \(x\neq1\).

\textbf{(c)} \(\dfrac{4}{x}+\dfrac{3}{x+2}=1\). Multiplicamos por \(x(x+2)\):


\[
4(x+2)+3x = x(x+2) \Rightarrow 4x+8+3x = x^2+2x
\]




\[
7x+8 = x^2+2x \Rightarrow x^2-5x-8=0.
\]


Resolviendo con la fórmula cuadrática:


\[
x=\frac{5\pm\sqrt{25+32}}{2}=\frac{5\pm\sqrt{57}}{2}.
\]


Comprobar que no anulan denominadores.
\end{solution}

\begin{solution}[Ejercicio 2]
\textbf{(a)} \(x^2-5x+6=(x-2)(x-3)=0\Rightarrow x=2,3.\)

\textbf{(b)} Discriminante \(\Delta=k^2-4\). Raíces reales y distintas si \(\Delta>0\Rightarrow k^2>4\Rightarrow k>2\) o \(k<-2\).

\textbf{(c)} Usamos fórmula:


\[
x=\frac{m+1\pm\sqrt{(m+1)^2-4m}}{2}=\frac{m+1\pm\sqrt{m^2-2m+1}}{2}=\frac{m+1\pm|m-1|}{2}.
\]


Entonces:


\[
\begin{cases}
\text{Si } m\ge1: x=\dfrac{m+1\pm(m-1)}{2}\Rightarrow x\in\{m,1\}.\

\[4pt]
\text{Si } m<1: x=\dfrac{m+1\pm(1-m)}{2}\Rightarrow x\in\{1,m\}.
\end{cases}
\]


(Se obtiene el mismo conjunto: \(x=1\) y \(x=m\).)
\end{solution}

\begin{solution}[Ejercicio 3]
Sustituimos \(x=y+1\) en la primera ecuación o resolvemos por sustitución:


\[
x-y=1\Rightarrow x=y+1.
\]


Sustituir en \(x^2+y^2=25\):


\[
(y+1)^2+y^2=25 \Rightarrow 2y^2+2y+1-25=0 \Rightarrow 2y^2+2y-24=0
\]




\[
y^2+y-12=0 \Rightarrow (y+4)(y-3)=0 \Rightarrow y=-4 \text{ o } y=3.
\]


Entonces \(x=y+1\) da \(x=-3\) o \(x=4\). Soluciones: \((-3,-4)\) y \((4,3)\).
\end{solution}

\begin{solution}[Ejercicio 4]
\textbf{(a)} \(2^{x+1}=8=2^3\Rightarrow x+1=3\Rightarrow x=2.\)

\textbf{(b)} \(3^{2x-1}=27=3^3\Rightarrow 2x-1=3\Rightarrow x=2.\)

\textbf{(c)} \(\log_2\frac{x+3}{x-1}=2\Rightarrow \frac{x+3}{x-1}=2^2=4\).


\[
x+3=4x-4\Rightarrow -3x=-7\Rightarrow x=\frac{7}{3}.
\]


Verificar dominio: \(x>1\) cumple.
\end{solution}

\begin{solution}[Ejercicio 5]
Primero restricciones: \(x\neq\pm1\).



\[
\frac{x^2-4}{x^2-1}=\frac{3x-2}{x+1}.
\]


Factorizamos \(x^2-4=(x-2)(x+2)\) y \(x^2-1=(x-1)(x+1)\):


\[
\frac{(x-2)(x+2)}{(x-1)(x+1)}=\frac{3x-2}{x+1}.
\]


Multiplicamos por \((x-1)(x+1)\):


\[
(x-2)(x+2)= (3x-2)(x-1).
\]


Expandimos:


\[
x^2-4 = 3x^2 -3x -2x +2 = 3x^2 -5x +2.
\]


Llevamos todo a un lado:


\[
0 = 2x^2 -5x +6.
\]


Resolvemos:


\[
x = \frac{5\pm\sqrt{25-48}}{4} = \frac{5\pm\sqrt{-23}}{4}.
\]


No hay soluciones reales. Concluimos que no existen soluciones reales; verificar que no se ha perdido ninguna por restricciones.
\end{solution}

\begin{solution}[Ejercicio 6]
\textbf{(a)} \(y=kx^2\). Dado \(y(2)=8\Rightarrow 8=k\cdot 4\Rightarrow k=2\). Entonces \(y(5)=2\cdot 25=50.\)

\textbf{(b)} \(t=\dfrac{c}{v}\). Dado \(t(4)=3\Rightarrow 3=\dfrac{c}{4}\Rightarrow c=12\). Entonces \(t(12)=\dfrac{12}{12}=1.\)

\textbf{(c)} \(I = k\frac{P}{d^2}\). Dado \(50 = k\frac{100}{2^2}=k\frac{100}{4}=25k\Rightarrow k=2.\) Para \(P=150,d=3\):


\[
I = 2\cdot\frac{150}{9} = \frac{300}{9} = \frac{100}{3}\approx 33.33.
\]


\end{solution}

\begin{solution}[Ejercicio 7]
\textbf{(a)} Tasa primera manguera \(= \tfrac{1}{3}\) tanque/h, segunda \(=\tfrac{1}{5}\). Juntas:


\[
\frac{1}{3}+\frac{1}{5}=\frac{5+3}{15}=\frac{8}{15}\ \text{tanques/h}.
\]


Tiempo total \(= \dfrac{1}{8/15}=\dfrac{15}{8}=1.875\ \text{h}=1\ \text{h }52.5\ \text{min}.\)

\textbf{(b)} Relación \(A:B=3:5\). Total partes \(=8\). Para 40 kg:


\[
A=\frac{3}{8}\cdot 40 = 15\ \text{kg},\quad B=\frac{5}{8}\cdot 40 = 25\ \text{kg}.
\]



\textbf{(c)} Velocidad \(v(t)=s'(t)=-9.8t\). Queremos \(v=-19.6\):


\[
-9.8t=-19.6\Rightarrow t=2\ \text{s}.
\]


Posición \(s(2)=100-4.9(4)=100-19.6=80.4\ \text{m}.
\]


\end{solution}

\bigskip
\section*{Consejos para el docente}
\begin{itemize}
  \item Añade variaciones numéricas para generar hojas de ejercicios distintas.
  \item Para evaluación automática, exporta las soluciones a un archivo separado y usa comparadores simbólicos.
  \item Incluye diagramas con TikZ para problemas geométricos relacionados con proporciones.
\end{itemize}

\end{document}

